\chapter*{Prefácio e Agradecimentos}

Este é um texto preliminar destinado a uma disciplina sobre sistemas operacionais.  
Ele explica os principais conceitos de sistemas operacionais por meio do estudo de um kernel de exemplo, chamado xv6.  
O xv6 é inspirado no Unix Version 6 (v6) de Dennis Ritchie e Ken Thompson~\cite{unix}.  
O xv6 segue de forma livre a estrutura e o estilo do v6, mas é implementado em ANSI C~\cite{kernighan} para uma arquitetura RISC-V multicore~\cite{riscv}.

Este texto deve ser lido em conjunto com o código-fonte do xv6,  
uma abordagem inspirada no comentário de John Lions sobre o UNIX 6ª Edição~\cite{lions};  
o texto possui hiperlinks para o código-fonte em  
\url{https://github.com/mit-pdos/xv6-riscv}.  
Consulte também  
\url{https://pdos.csail.mit.edu/6.1810}  
para recursos online adicionais sobre o v6 e o xv6, incluindo diversos exercícios práticos utilizando o xv6.

Utilizamos este texto nos cursos 6.828 e 6.1810, disciplinas de sistemas operacionais no MIT.  
Agradecemos ao corpo docente, assistentes de ensino e estudantes dessas turmas, que direta ou indiretamente contribuíram com o xv6.  
Em particular, gostaríamos de agradecer a Adam Belay, Austin Clements e Nickolai Zeldovich.  
Finalmente, agradecemos às pessoas que nos enviaram relatórios de erros no texto ou sugestões de melhorias:  
Abutalib Aghayev, Sebastian Boehm, brandb97, Anton Burtsev, Raphael Carvalho, Tej Chajed, Brendan Davidson, Rasit Eskicioglu, Color Fuzzy, Wojciech Gac, Giuseppe, Tao Guo, Haibo Hao, Naoki Hayama, Chris Henderson, Robert Hilderman, Eden Hochbaum, Wolfgang Keller, Paweł Kraszewski, Henry Laih, Jin Li, Austin Liew, lyazj@github.com, Pavan Maddamsetti, Jacek Masiulaniec, Michael McConville, m3hm00d, miguelgvieira, Mark Morrissey, Muhammed Mourad, Harry Pan, Harry Porter, Siyuan Qian, Zhefeng Qiao, Askar Safin, Salman Shah, Huang Sha, Vikram Shenoy, Adeodato Simó, Ruslan Savchenko, Pawel Szczurko, Warren Toomey, tyfkda, tzerbib, Vanush Vaswani, Xi Wang, Zou Chang Wei, Sam Whitlock, Qiongsi Wu, LucyShawYang, ykf1114@gmail.com e Meng Zhou.

Se você encontrar erros ou tiver sugestões de melhorias, por favor envie um e-mail para  
Frans Kaashoek e Robert Morris (kaashoek, rtm@csail.mit.edu).
